\vspace{-0.2cm}

\section{Alternative Views}
% A popular alternative view, especially in industry, is that there are very limited further innovations required to achieve an AI software engineer. In the past few years, we have seen emergent capabilities come from scaling \citep{wei2022emergent}, and this view may posit that the majority of the tasks and challenges we pose in this work can be resolved by scaling to larger models and training with more data.
An alternative view, particularly within the industry, suggests that limited further innovation is needed to achieve AI software engineers. They believe that scaling up larger models and more data is sufficient for creating AI engineers capable of contributing effectively. This is supported by the significant progress in recent years, for example the SWE-bench score increasing from an initial 2\% to now 73\%. However, we believe this view is too optimistic. As we have outlined in the paper, there are many fundamental capabilities in AI for SWE that today's models lack. As we argued in Sec. \ref{sec:main-paths}, improving these aspects requires fundamentally new forms of data, training environments, and inference-time approaches. The alternative perspective, while not unreasonable given recent progress, should not overshadow the need for more efforts in what we outline. As a research field, AI for code still has plenty to offer and overcome.